\documentclass[11pt]{article}
\usepackage[margin=1.1in]{geometry}
\usepackage{amsmath,amssymb,amsthm}
\usepackage{booktabs}
\newtheorem{theorem}{Theorem}
\newtheorem{lemma}{Lemma}
\newtheorem{corollary}{Corollary}
\newtheorem{proposition}{Proposition}
\theoremstyle{remark}
\newtheorem{remark}{Remark}
\newcommand{\odd}{\operatorname{odd}}
\newcommand{\pos}[1]{(#1)^{+}}
\title{Circuit codes of dimension near $3k/2$:\\ an upper bound by window counting, and counterexamples\\ to a conjecture of Byrnes}
\author{Claude (an AI model made by Anthropic), working autonomously}
\date{Draft of 12 November 2026 --- not reviewed by a human mathematician;\\ the main statements are machine-checked in Lean (Section~\ref{sec:formal})}
\begin{document}
\maketitle
\begin{abstract}
A circuit code of spread $k$ is a cycle in the $n$-cube $Q_n$ in which any two vertices at Hamming distance less than
$k$ are at the same distance along the cycle; $K(n,k)$ is the maximum length. Write $n=(3k+l+1)/2$. Exact values
were known for $l\le 3$ and, for $k\ge 9$, $l=4$ (Singleton, 1966; Douglas, 1969). Byrnes constructed codes of length $4k+2l$ for $k\ge 2l+1$ ($k$ odd) and
$k\ge 2l-2$ ($k$ even) and conjectured that they are optimal. We show:
(i) by explicit codes, the conjecture fails for more than $60$ pairs $(n,k)$ on the lines $l=6,\dots,17$;
(ii) by an elementary double count over windows of the transition sequence,
$K(n,k)\le 4k+2l$ whenever $k>(l^2-1)/2$ ($l$ odd) or $k>(l-2)(l+1)/2$ ($l$ even), so the conjecture is true
for all but finitely many $k$ on every line;
(iii) a code of length $4k+2l+2$ has a rigid shape (a balanced two-letter ``type word''), which reduces its existence to a
small satisfiability problem. With (iii) the lines $l=5,\dots,9$ are settled completely: the conjecture fails exactly for
$k=13$ ($l=6$), $k=12,\dots,18$ ($l=7$), $k=17,\dots,25$ ($l=8$) and for every admissible $k=16,\dots,40$ on the line $l=9$,
where the bound (ii) is therefore sharp;
(iv) an explicit family of codes of length $4k+2l+2$ at $k=(l^2-1)/2$ for every $l\equiv 1\pmod 4$, $l\ge 9$: the bound (ii) is
sharp for infinitely many $l$, the conjecture fails on every line $l\ge 6$, and the largest $k$ for which it fails on the line $l$ is $l^2/2-O(l)$. We obtain $K(n,k)$ exactly on these lines for $k\ge 20$ ($l=7$) and $k\ge 19$ ($l=8$), including $K(33,19)=96=4k+2l+4$,
and on the line $l=9$ for $k\ge 26$.
(v) (Section~\ref{sec:last}, added last) for $l\equiv 0,3\pmod 4$ no code of length $4k+2l+2$ exists for $k>(l^2-l-6)/2$, and explicit codes
exist at $k=(l^2-l-6)/2$: the last $k$ at which the conjecture fails is $(l^2-1)/2$ for $l\equiv 1$ and $(l^2-l-6)/2$ for $l\equiv 0,3 \pmod 4$;
for $l\equiv 2\pmod 4$ it is $(l^2-l-8)/2$ or $(l^2-l-4)/2$ (the latter for all $l\le 38$).
(vi) (Section~\ref{sec:sym}) Byrnes' bound for symmetric codes with a bit run $k+l$ holds for all $k$, with a short proof, so
Conjecture~4.2 is equivalent to $K(n,k)=4k+2l$; it fails on every line $l\ge 6$, by an unbounded amount.
The main upper bounds, the families for $l\equiv 1\pmod 4$, and (vi) are checked in Lean~4 (Section~\ref{sec:formal}).
\end{abstract}

\section{Definitions and results}
Let $C=(v_0,\dots,v_{N-1})$ be a cycle in $Q_n$ (indices mod $N$), $d_H$ the Hamming distance and
$d_C(i,j)=\min(|i-j|,N-|i-j|)$. $C$ is a \emph{circuit code of spread $k$} (a $k$-coil) if
\[ d_H(v_i,v_j)<k \;\Longrightarrow\; d_C(i,j)=d_H(v_i,v_j)\qquad\text{for all } i,j. \]
$K(n,k)$ denotes the maximum length $N$ of such a code. Throughout, $k\ge 2$,
\[ l=2n-3k-1,\qquad\text{i.e. } n=\tfrac{3k+l+1}{2},\qquad K=k+1, \]
so $k$ and $l$ have opposite parity. $K(n,k)=2n$ for $l\le 1$ (Singleton), and Douglas determined $K(n,k)$ for $l=2,3$ and for $l=4$, $k\ge 9$
(the value is $4k+2l$).
Byrnes \cite{B2} constructs, for every $l\ge 2$ and every $k\ge 2l+1$ ($k$ odd) or $k\ge 2l-2$ ($k$ even), a code of
length $4k+2l$, proves that this is the maximum among codes with a certain symmetry and bit run $k+l$, and conjectures
(Conjecture~4.2 of \cite{B2}) that $K(n,k)=4k+2l$ in this whole range.

\begin{theorem}\label{main}
Let $l\ge 2$, $n=(3k+l+1)/2$, and suppose
\[ k>\frac{l^2-1}{2}\ \text{ and }\ k>3l-3\quad(l\text{ odd}),\qquad\text{or}\qquad k>\frac{(l-2)(l+1)}{2}\quad(l\text{ even}). \]
Then every spread-$k$ circuit code in $Q_n$ has length at most $4k+2l$. (For odd $l\geq 5$ the condition $k>3l-3$ is implied by the first.)
\end{theorem}

With Byrnes' construction, $K(n,k)=4k+2l$ in the range of Theorem~\ref{main} when $l\ge 5$.
The largest $k$ \emph{not} covered is $12,13,24,27,40,43,60,65$ for $l=5,\dots,12$.

\begin{theorem}\label{ce}
$K(n,k)\ge 4k+2l+2$ for the pairs of Table~\ref{tab}, all of which lie in the range of Byrnes' conjecture; and $K(33,19)\ge 96$, $K(32,18)\ge 94$ (excess $4$).
\end{theorem}

\begin{table}[h]\centering
\begin{tabular}{cp{0.84\textwidth}}\toprule
$l$ & $(n,k)$ with a code of length $4k+2l+2$ \quad($^*$: also $4k+2l+4$) \\ \midrule
6 & $(23,13)$ \\
7 & $(22,12)$, $(25,14)$, $(28,16)$, $(31,18)$ \\
8 & $(30,17)$, $(33,19)^*$, $(36,21)$, $(39,23)$, $(42,25)$ \\
9 & $(29,16)$, $(32,18)^*$, $(35,20)$, $(38,22)$, $(41,24)$, $(44,26)$, $(47,28)$, $(50,30)$, $(53,32)$, $(56,34)$, $(59,36)$, $(62,38)$, $(65,40)$ \\
10 & $(37,21)$, $(40,23)$, $(43,25)^*$, $(46,27)$, $(49,29)$, $(52,31)$, $(55,33)$, $(58,35)$, $(61,37)$, $(64,39)$ \\
11 & $(36,20)$, $(39,22)$, $(42,24)$, $(45,26)$, $(48,28)$ \\
12 & $(44,25)$, $(47,27)$ \\
13 & $(43,24)$, $(46,26)$ \\ \bottomrule
\end{tabular}
\caption{Counterexamples to Conjecture 4.2 of \cite{B2} (state of 10 Nov 2026, 13:25; lines $l\ge 11$ incomplete, further pairs on $l=11,\dots,17$ follow by contraction: see \texttt{listce.py}). The codes are in the files
\texttt{codes/sat\_*}, \texttt{pump\_*}, \texttt{con\_*}, \texttt{skel\_*};
each was validated by three independent programs directly from the definition.}\label{tab}
\end{table}

The first entry of each line is the smallest $k$ admitted by the conjecture ($k=2l+1$ or $2l-2$). An explicit
$(23,13)$ code of length $66$ is given in the Appendix.

\begin{corollary}\label{lines}
(a) $K\bigl(\tfrac{3k+6}{2},k\bigr)=4k+10$ for every even $k\ge 8$.\quad
(b) $K(23,13)=66$, and $K\bigl(\tfrac{3k+7}{2},k\bigr)=4k+12$ for every odd $k\ge 15$.
\end{corollary}
Here $k\ge 14$ in (a) and $k\ge 15$ in (b) follow from Theorem~\ref{main}; the cases $K(15,8)=42$, $K(18,10)=50$, $K(21,12)=58$
and the upper bound $K(23,13)\le 66$ are computer proofs (SAT, with DRAT certificates checked by \texttt{drat-trim}; see Section~\ref{comp}). So the conjecture is true on the line $l=5$,
and on the line $l=6$ it fails exactly at its first case $k=13$, which is also the last $k$ not covered by Theorem~\ref{main}.

\begin{theorem}\label{l78}
(a) $K(34,20)=94$, $K(37,22)=102$, $K(40,24)=110$. Hence on the line $l=7$: $K\bigl(\tfrac{3k+8}{2},k\bigr)=4k+14$ for every even
$k\ge 20$, while $K>4k+14$ for $k=12,14,16,18$.\\
(b) $K(33,19)=96$, $K(36,21)=102$, $K(39,23)=110$, $K(42,25)=118$, $K(45,27)=124$. Hence on the line $l=8$:
$K\bigl(\tfrac{3k+9}{2},k\bigr)=4k+16$ for every odd $k\ge 27$, while $K>4k+16$ for $k=17,\dots,25$.\\
(c) On the line $l=9$: $K\bigl(\tfrac{3k+10}{2},k\bigr)=4k+20$ for every even $k$ with $26\le k\le 40$ (and $K\ge 4k+20$ for $16\le k\le 24$), while
$K=4k+18$ for $k\ge 42$. In particular $K(65,40)=180$, the case $2k+1=l^2$ excluded from Theorem~\ref{main}: the hypothesis of
Theorem~\ref{main} cannot be weakened for $l=9$ (nor for $l=6$).
\end{theorem}
The upper bounds in Theorem~\ref{l78} combine Lemmas~\ref{L1} and \ref{L2} with Proposition~\ref{shape} and a computer search
(Section~\ref{sec:shape}); $K(40,24)=110$ is proved by hand there.

\section{Windows}
Let $T=(T_0,\dots,T_{N-1})$ be the cyclic transition sequence: $T_i$ is the coordinate in which $v_i$ and $v_{i+1}$
differ. A \emph{window} $W$ of length $L<N$ is a set of $L$ cyclically consecutive positions $t,\dots,t+L-1$;
it joins $v_t$ to $v_{t+L}$. Let $\odd(W)$ be the number of coordinates occurring an odd number of times in $W$, so
$\odd(W)=d_H(v_t,v_{t+L})$ and $\odd(W)\equiv L \pmod 2$.

\medskip\noindent\textbf{(F1)} If $k\le L\le N-k$ then $\odd(W)\ge k$, hence $\odd(W)\ge k+\bigl((L-k)\bmod 2\bigr)$.

Indeed $d_C(t,t+L)=\min(L,N-L)\ge k$, and $d_H<k$ would force $d_C=d_H<k$.

\medskip\noindent\textbf{(F2)} If $N\ge 2k+1$, any $K=k+1$ consecutive transitions are distinct.

For $j\le k-1$ a window of length $j$ joins vertices with $d_C=j$ and $d_H\le j$; $d_H<j$ would give $d_C=d_H<j$.
So the $j$ transitions are distinct. For $j=k$ and $j=k+1$ use (F1): $\odd\ge k$ and $\odd\equiv j$.

\medskip\noindent\textbf{(F3)} For a position $i$ let $i^+$ be the next position (cyclically) carrying the same coordinate,
$g_i=i^+-i\ge K$ the \emph{gap}, and $e_i=g_i-K\ge 0$. A coordinate used $u$ times has $u$ gaps summing to $N$ ($u$ is even,
since the walk closes; $u\geq 2$ if the coordinate is used at all). Hence
$\sum_i g_i = N\cdot\#\{\text{used coordinates}\}\le Nn$. The occurrences $i,i^+,\dots,i^{+(j-1)}$ of one coordinate
(with $j\le u$, total span $G=g_i+\dots+g_{i^{+(j-2)}}<N$) lie together in exactly $\pos{L-G}$ of the $N$ windows of length $L$,
where $\pos{x}=\max(x,0)$. A window of length $L\le jK$ contains each coordinate at most $j$ times.

\medskip We use two elementary facts about $\pos{\cdot}$: $\pos{a+b}\le\pos a+\pos b$; and if $c\ge a$ and $c\ge b$ then
\begin{equation}\label{el} \pos{a+b-c}\ \ge\ \pos a+\pos b-\pos c . \end{equation}
(If $a,b\ge 0$ then $c\ge0$ and the right side is $a+b-c$. If $a<0\le b$ the right side is $b-\pos c= b-c\le 0$; similarly
for $b<0\le a$. If $a,b<0$ the right side is $-\pos c\le 0$.)

\begin{lemma}[chord inequality]\label{L1}
Let $C$ be a spread-$k$ circuit code of length $N$ in $Q_n$ with $2k+2\le N<6K$, and let $L$ be an integer with
$\max(k,N/4)\le L\le\min(2K,N-k)$. Then
\[ \Bigl(\frac N2-n\Bigr)(4L-N)\ \le\ N\Bigl\lfloor\frac{L-k}{2}\Bigr\rfloor . \]
\end{lemma}
\begin{proof}
Since $L\le 2K$, every coordinate occurs at most twice in a window $W$ of length $L$, so $\odd(W)=L-2P(W)$ where
$P(W)$ is the number of coordinates occurring twice, i.e.\ the number of pairs $(i,i^+)$ inside $W$ (``chords'').
By (F1), $P(W)\le\lfloor (L-k)/2\rfloor$. By (F3), $\sum_W P(W)=\sum_i\pos{L-g_i}$.
Since $N<6K$, by (F2) every coordinate is used $0$, $2$ or $4$ times. If $n_1$ coordinates are used twice and $n_2$ four times,
then $n_1+n_2\le n$ and $2n_1+4n_2=N$, so $n_2\ge N/2-n$. The four gaps of a coordinate used four times sum to $N$, so
they contribute $\sum\pos{L-g}\ge 4L-N\ge 0$; all other terms are non-negative. Therefore
\[ \Bigl(\frac N2-n\Bigr)(4L-N)\le n_2(4L-N)\le\sum_W P(W)\le N\Bigl\lfloor\frac{L-k}{2}\Bigr\rfloor. \qedhere\]
\end{proof}

\begin{lemma}[no long codes]\label{L2}
Let $C$ be a spread-$k$ circuit code of length $N\ge 5k+4$ in $Q_n$. Then
\[ k+(k\bmod 2)\ \le\ 6n-12K+6\Bigl\lfloor\frac{k+2}{2}\Bigr\rfloor . \]
With $n=(3k+l+1)/2$ the right side equals $3l-3$ for $k$ even and $3l-6$ for $k$ odd. Consequently a code with
$k>3l-3$ ($k$ even) or $k>3l-7$ ($k$ odd) has length $N\le 5k+3$.
\end{lemma}
\begin{proof}
$N\ge 5k+4$ guarantees $k\le L\le N-k$ for $L=2K,3K,4K$. Put $A=\lfloor (k+2)/2\rfloor$,
$x_i=\pos{K-e_i}$ and $b_i=\pos{K-e_i-e_{i^+}}$.

\emph{(a) $L=2K$.} As in Lemma~\ref{L1}, $P(W)\le\lfloor(2K-k)/2\rfloor=A$ and
$\sum_WP(W)=\sum_i\pos{2K-g_i}=\sum_ix_i$, so
\begin{equation}\label{e1}\sum_i x_i\le NA.\end{equation}

\emph{(b) $L=3K$.} Each coordinate occurs at most three times in $W$; let $c_j$ be the number occurring exactly $j$ times.
Then $3K=c_1+2c_2+3c_3$, $\odd(W)=c_1+c_3$ and $c_1+c_2+c_3\le n$. So $c_2+c_3=(3K-\odd)/2$ and
$c_1+c_2+c_3=(3K+\odd)/2-c_3\le n$, whence, using $\odd\ge k$ and $\odd\equiv 3K$,
$c_3\ge\lceil(3K+k)/2\rceil-n=2K-n$. As $c_3(W)$ is the number of triples $(i,i^+,i^{++})$ inside $W$, (F3) gives
\begin{equation}\label{e2}\sum_i b_i=\sum_i\pos{3K-g_i-g_{i^+}}=\sum_Wc_3(W)\ge N(2K-n).\end{equation}
(For a coordinate used twice, $g_i+g_{i^+}=N>3K$ and the term vanishes, as it should.)

\emph{(c) $L=4K$.} Each coordinate occurs at most four times. With $c_j$ as before,
$4K=c_1+2c_2+3c_3+4c_4$ and
\[ \odd(W)=c_1+c_3=4K-2\,\mathrm{ch}(W)+2\,t(W)-2\,q(W), \]
where $\mathrm{ch}=c_2+2c_3+3c_4$ is the number of chords $(i,i^+)$ inside $W$, $t=c_3+2c_4$ the number of triples
$(i,i^+,i^{++})$ inside $W$ and $q=c_4$ the number of quadruples $(i,i^+,i^{++},i^{+++})$ inside $W$. Summing over the $N$ windows with (F3):
\begin{align*}
\sum_W\mathrm{ch}&=\sum_i\pos{4K-g_i}\ \ge\ \sum_i(4K-g_i)\ \ge\ N(4K-n);\\
\sum_Wt&=\sum_i\pos{(K-e_i)+(K-e_{i^+})}\ \le\ \sum_i(x_i+x_{i^+})=2\sum_ix_i\ \le\ 2NA;\\
\sum_Wq&=\sum_i\pos{K-e_i-e_{i^+}-e_{i^{++}}}\ \ge\ \sum_i\bigl(b_i+b_{i^+}-x_{i^+}\bigr)\\ &=2\sum_ib_i-\sum_ix_i\ \ge\ 2N(2K-n)-NA.
\end{align*}
The second line uses that $i\mapsto i^+$ is a bijection and \eqref{e1}; the third uses \eqref{el} with
$a=K-e_i-e_{i^+}$, $b=K-e_{i^+}-e_{i^{++}}$, $c=K-e_{i^+}$, then \eqref{e1} and \eqref{e2}.
(A coordinate used twice has no triples or quadruples; its terms in the sums for $t$ and $q$ are $0$, since the spans are
$\ge N>4K$, and the bounds remain valid for them: the upper bound $x_i+x_{i^+}$ is $\ge 0$, and the lower bound
$b_i+b_{i^+}-x_{i^+}=-x_{i^+}$ is $\le 0$ because $b_i=b_{i^+}=\pos{3K-N}=0$.) Hence
\[ \sum_W\odd(W)\le N\bigl[4K-2(4K-n)+4A-4(2K-n)+2A\bigr]=N\,[\,6n-12K+6A\,], \]
while $\odd(W)\ge k+(k\bmod 2)$ for every $W$ by (F1) with $L=4K$ even.
For $k$ even $6A=3K+3$ and $6n=9K+3l-6$, giving $3l-3$; for $k$ odd $6A=3K$, giving $3l-6$.
\end{proof}

\begin{proof}[Proof of Theorem~\ref{main}]
Suppose a code has length $N=4k+2l+e$ with $e\ge 2$ ($N$ is even).

If $N\ge 5k+4$, Lemma~\ref{L2} is violated: for $k$ even $k>3l-3$ holds by hypothesis; for $k$ odd,
$k>(l-2)(l+1)/2\ge 3l-7$ because $(l-3)(l-4)\ge 0$.

Otherwise $2\le e\le k+4-2l$ and $N<6K$, so Lemma~\ref{L1} applies; note $N/2-n=(k+l-1+e)/2$.

\emph{$k$ even.} Take $L=2k+1$ (then $\lfloor(L-k)/2\rfloor=k/2$ and $4L-N=4k+4-2l-e>0$). Lemma~\ref{L1} says $f(e)\le 0$ for
\[ f(e)=(k+l-1+e)(4k+4-2l-e)-(4k+2l+e)k=-e^2+(2k+5-3l)e-(2l^2-6l+4). \]
But $f$ is concave, $f(2)=4k+2-2l^2>0$ since $k>(l^2-1)/2$, and $f(k+4-2l)=k(k+5-3l)>0$ since $k>3l-3$.
So $f>0$ on $[2,k+4-2l]$, a contradiction.

\emph{$k$ odd.} Take $L=2k+2$ (then $\lfloor(L-k)/2\rfloor=(k+1)/2$ and $4L-N=4k+8-2l-e>0$). Lemma~\ref{L1} says $f(e)\le 0$ for
\[ f(e)=(k+l-1+e)(4k+8-2l-e)-(4k+2l+e)(k+1)=-e^2+(2k+8-3l)e-(2l^2-8l+8). \]
Now $f(2)=4k+4+2l-2l^2>0$ since $k>(l-2)(l+1)/2$, and $f(k+4-2l)=k^2+(8-3l)k+8-4l$. For $l=2$ this is $k(k+2)>0$.
For even $l\ge 4$, $k\ge(l-2)(l+1)/2+1$ gives $k+8-3l\ge(l^2-7l+16)/2\ge 2$, so
$f(k+4-2l)\ge 2k-(4l-8)\ge l^2-5l+8>0$. Again $f>0$ on the whole range, a contradiction.
(The window lengths are admissible: $N/4\le L\le 2K$ and $L\le N-k$ hold because $2\le e\le k+4-2l$ and $l\ge 2$.)
\end{proof}

\begin{remark}
For $l=2,3,4$ Theorem~\ref{main} gives $K(n,k)\le 4k+2l$ for odd $k\geq 3$, even $k\ge 8$, odd $k\ge 7$ respectively, in agreement
with the classical exact values.
\end{remark}
\begin{remark}[structure at equality]
If $f(2)=0$ (e.g.\ $(n,k)=(40,24)$, $l=7$) a code of length $4k+2l+2$ must have: every window of length $2k+1$ at
Hamming distance exactly $k+1$\,; every coordinate used; exactly $k$ coordinates used twice, with both gaps $\ge 2k+1$; and all
gaps of the $n-k$ four-times coordinates $\le 2k+1$. All sixteen codes of Theorem~\ref{ce}, and all excess-$2$ codes we have seen, have
exactly $k$ twice-used and $n-k$ four-times-used coordinates.
\end{remark}

\section{The shape of a code just above $4k+2l$}\label{sec:shape}
Let $C$ be a spread-$k$ circuit code of even length $N$ with $2K\le N<6K$, put $m=N/2-n$, $L_0=2k+1$ ($k$ even) or $2k+2$ ($k$ odd),
$c_0=\lfloor (L_0-k)/2\rfloor$, and assume $N/4\le L_0\le N-k$ and
\begin{equation}\label{H} (m+1)(4L_0-N)>Nc_0 . \end{equation}
For $N=4k+2l+2$ we have $m=n-k$, and \eqref{H} holds in every case used below.
\begin{proposition}\label{shape}
Under \eqref{H}: exactly $m$ coordinates are used four times (``$F$ coordinates''), $n-m$ are used twice (``dots''), none is unused,
and every gap of an $F$ coordinate is at most $g_{\max}=N-3K$. For a window $W$ of length $L$ let $F(W)$ be the number of its
positions carrying $F$ coordinates, $s(W)$ the number of dots occurring twice in $W$ and $t(W)$ the number of $F$ coordinates
occurring three times in $W$. If $g_{\max}\le L\le 3K$ then
\begin{equation}\label{id} \odd(W)=L+2m-2F(W)-2s(W)+2t(W), \end{equation}
and therefore, if also $k\le L\le N-k$,
\begin{equation}\label{bal} F(W)\ \le\ \Bigl\lfloor\frac{L+2m-k}{2}\Bigr\rfloor+t(W)-s(W). \end{equation}
Here $t(W)=0$ for $L\le 2K$; and for $L=2K+1+j$, $W=[a,a+L)$, $t(W)$ is at most the number of $F$ positions in each of
$[a,a+j]$, $[a+2K,a+2K+j]$ and $[a+K,a+L-1-K]$.
\end{proposition}
\begin{proof}
As in Lemma~\ref{L1}, uses are $0,2,4$, the number $n_2$ of four-times coordinates is at least $m$ and satisfies
$n_2(4L_0-N)\le Nc_0$; by \eqref{H} $n_2=m$, so $n_1=(N-4m)/2=n-m$ and all $n$ coordinates are used. Four gaps $\ge K$ summing to $N$
are each $\le N-3K$. If $L\ge g_{\max}$ every $F$ coordinate occurs in $W$, and at most three times if $L\le 3K$; with $a_j$ of them
occurring $j$ times, $a_1+a_2+a_3=m$ and $F(W)=a_1+2a_2+3a_3$ give $a_1+a_3=2m-F(W)+2a_3$; the dots contribute $L-F(W)-2s(W)$.
This is \eqref{id}; \eqref{bal} follows from (F1). Three occurrences $p<p'<p''$ of a coordinate in $W$ have gaps $\ge K$, so
$p\le a+j$, $p''\ge a+2K$, $a+K\le p'\le a+L-1-K$, and distinct coordinates use distinct positions.
\end{proof}
Summing \eqref{bal} over all windows of length $L_0$ gives Lemma~\ref{L1} again; but \eqref{bal} holds window by window and for many
lengths, and says that the cyclic two-letter \emph{type word} ($F$ or dot at each position) is balanced.

\paragraph{Example: $K(40,24)=110$ by hand.} Suppose $N=112$: $K=25$, $m=16$ ($64$ $F$ positions with gaps in $[25,37]$), $24$ dots.
By \eqref{bal} every window of length $49$ has $F(W)\le 28$, and the average of $F(W)$ is $49\cdot 64/112=28$; so $F(W)=28$ always,
the type word is invariant under the shift by $49$, hence has period $\gcd(49,112)=7$ with four $F$ and three dots per period.
A window of length $47$ has $F(W)\le 27$; it is a $49$-window without its two ends $t$, $t+48\equiv t-1 \pmod 7$, so these are not
both dots: no two cyclically adjacent dots, and the period is $\cdot FF\cdot F\cdot F$. Let $s,s+1$ be the residues of the two
adjacent $F$'s; then $s+4$ is a dot residue. A window $[p,p+50]$ of length $51$ with $p\equiv s$ has $F(W)=28+2=30>29$, so by
\eqref{bal} $t(W)\ge1$, which needs an $F$ at $p+25\equiv s+4$: contradiction. Lengths $\ge 114$ are excluded by Lemmas~\ref{L1} and \ref{L2}. \qed

\paragraph{The general case.} Relax the dots to anonymous positions. A code as in Proposition~\ref{shape} yields an
\emph{$F$-skeleton}: an assignment of $m$ coordinates and $n-m$ dot marks to the $N$ positions such that each coordinate is used an
even number of times, never twice within $K$ consecutive positions, at least once in every $g_{\max}$ consecutive positions, and
for every window $W$ with $k\le |W|\le N-k$
\[ D(W)+\odd_F(W)\ \ge\ k+\bigl((|W|-k)\bmod 2\bigr), \]
where $D(W)$ is the number of dots in $W$ and $\odd_F(W)$ the number of $F$ coordinates occurring an odd number of times
(since $\odd(W)=D(W)-2s(W)+\odd_F(W)$). Existence of an $F$-skeleton is a satisfiability problem with only $m$ coordinates;
if it is unsatisfiable there is no code of length $N$. For $(n,k,N)=(37,22,104)$, $(40,24,112)$, $(34,20,96)$, $(45,27,126)$ it is
unsatisfiable (kissat: $51$, $38$, $134$, $221$ seconds; DRAT proofs by cadical verified with drat-trim), whereas a direct encoding of
the definition did not terminate in $3.5$ hours. When it is satisfiable, a second small SAT problem assigns the dots; this produced
the codes for $(36,21)$, $(39,23)$, $(42,25)$, all pairs with $24\le k\le 40$ on the line $l=9$, $(43,25),\dots,(55,33)$ on $l=10$, and the
$(33,19)$ code of length $96$. For $(33,19)$ the length $98$ is excluded by Lemma~\ref{L1} and larger ones by Lemma~\ref{L2}.
In the whole region $k>3l-3$ ($k$ even), $k>3l-7$ ($k$ odd), $7\le l\le 12$, every length $4k+2l<N\le 5k+3$ either violates
Lemma~\ref{L1} or satisfies \eqref{H}, so there $K(n,k)$ is determined by at most three such queries.

\paragraph{An explicit family at $2k+1=l^2$.} Let $l\equiv 1 \pmod 4$, $l\ge 9$, $k=(l^2-1)/2$, $K=k+1$, $N=2l(l+1)=4k+2l+2$ and
$n=(3k+l+1)/2=(3l-1)(l+1)/4$. Define a cyclic word $T$ on $\mathbb{Z}_N$ over $n$ symbols by
\begin{itemize}
\item[(A)] $T(lj)=T(lj+K)=T(lj+2K)=T(lj+3K+1)=A_j$ for $j=0,\dots,2l+1$ (these positions are exactly those with residues
$0$, $\tfrac{l+1}{2}$, $1$, $\tfrac{l+5}{2}$ modulo $l$);
\item[(B)] $T(t)=B_{t \bmod N/4}$ if $t \bmod l\in\{3,5,\dots,l-2\}\setminus\{\tfrac{l+1}{2},\tfrac{l+5}{2}\}$;
\item[(D)] $T(t)=D_{t\bmod N/2}$ if $t\bmod l\in\{2,4,\dots,l-1\}$.
\end{itemize}
This uses $2l+2$ symbols $A_j$, $(l-7)(l+1)/4$ symbols $B$ (together $(l+1)^2/4=n-k$, each four times) and $k$ symbols $D$ (each twice,
at antipodal positions); the type word is the one forced by the argument for $(40,24)$ above.
\begin{theorem}\label{fam}
For every $l\equiv 1\pmod 4$, $l\ge 9$, $T$ is the transition sequence of a spread-$k$ circuit code of length $4k+2l+2$ in $Q_n$. Consequently
$K(n,k)=4k+2l+2$ for $k=(l^2-1)/2$, $n=(3k+l+1)/2$, and the condition $k>(l^2-1)/2$ in Theorem~\ref{main} cannot be weakened for these $l$.
\end{theorem}
\begin{proof}[Proof for $l\ge 13$ (the case $l=9$, and all $l\le 161$, were checked by computer from the definition)]
Put $h=(l-1)/2$, $Q=N/4=l(l+1)/2\equiv 0 \pmod l$, so $K=Q-h$. The gaps are $K,K,K+1,g_4:=Q+3h-1$ for the $A_j$, $Q$ for the $B$'s and $2Q$ for the $D$'s:
all are $\ge K$ and every symbol is used an even number of times, so $T$ closes up and any $K$ consecutive transitions are distinct.
It remains to show $\odd(W)\ge k$ for every window $W=[t,t+L)$ with $K<L\le N/2$. Such a window contains no $D$ twice.
Let $c_R(t,x)$ be the number of $p\in[t,t+x)$ with $p \bmod l\in R$.

\emph{Step 1: $L=K+y$, $1\le y\le 2l-3$ (i.e.\ $L\le g_4$).} No symbol occurs three times, so $\odd(W)=L-2\,\mathrm{ch}(W)$ with
$\mathrm{ch}(W)=c_{\{0\}}(t,y)+c_{\{(l+1)/2\}}(t,y)+c_{\{1\}}(t,y-1)+c_{\mathcal B}(t,y-h)$ ($\mathcal B$ the residues of type (B)), and we need $2\,\mathrm{ch}\le y+1$.
Positions that start a pair have residue $0$ or an odd residue, so two adjacent positions both start pairs only for the residues $(0,1)$;
and the residue $(l+1)/2$ is at cyclic distance $\ge h$ from the residues $0$ and $1$.
If $y\le h-1$ only the residues $0,1,(l+1)/2$ count and at most the pair $(0,1)$ fits: $\mathrm{ch}\le 2$, and $\mathrm{ch}\le1$ for $y\le 2$.
If $y\ge h$, the last $h$ positions contribute at most $2$ and the first $y-h$ positions at most $\lceil (y-h-J)/2\rceil+J$, where $J\le 2$ is the
number of pairs $(0,1)$ among them; so $2\,\mathrm{ch}\le y-h+J+5\le y+1$ as $h\ge 6$.

\emph{Step 2: $g_4<L\le 2K$.} By the counting that gives \eqref{id} (it does not need $T$ to be a code; every $F$ symbol occurs once or twice in $W$) $\odd(W)\ge k$ iff $F(W)\le\lfloor (L+l+1)/2\rfloor$.
Write $L=lq+\rho$, $0\le\rho<l$. The type word has $(l+1)/2$ letters $F$ per period and at most $\lfloor\rho/2\rfloor+1$ in $\rho\ge1$ consecutive positions.
The inequality becomes $q\le l-1+(\rho\bmod 2)-((q+\rho)\bmod 2)$ for $\rho\ge 1$ and $q\le l+1-(q\bmod 2)$ for $\rho=0$; it holds for $q\le l-2$, for $q=l-1$ (even),
and for $q=l$, which occurs only with $\rho\in\{0,1\}$.

\emph{Step 3: $L=2K+z$, $1\le z\le l-1$.} A symbol occurs three times in $W$ exactly for $A_j$ with $lj\in[t,t+z)$ or $lj+K\in[t,t+z-1)$; let $\tau$ be their number.
By \eqref{id}, $\odd(W)\ge k$ iff $F(W)\le\lfloor (L+l+1)/2\rfloor+\tau$, i.e.\ $F'\le 1+\lfloor z/2\rfloor+\tau$, where $F'$ is the number of $F$ letters among the
residues $t,\dots,t+z$. If these do not contain $0$ followed by $1$, they alternate and $F'\le\lfloor z/2\rfloor+1$; otherwise $F'\le\lfloor (z-1)/2\rfloor+2$ and $\tau\ge1$.

The upper bound $K(n,k)\le 4k+2l+2$ is Lemma~\ref{L1} with $L=2k+1$ ($f(2)=0$, $f$ concave, $f(k+4-2l)>0$) and Lemma~\ref{L2}.
\end{proof}
\begin{corollary}
Conjecture 4.2 of \cite{B2} fails on every line $l\ge 6$. More precisely $k_{\max}(l)\ge \tfrac{(l-3)^2-1}{2}-3$ for $l\ge 9$, while $k_{\max}(l)\le\tfrac{l^2-1}{2}$ by Theorem~\ref{main}.
\end{corollary}
\begin{proof}
Deleting a twice-used coordinate (Deimer) from a code of length $4k+2l+2$ on the line $l$ gives a code of length $4(k-1)+2(l+1)+2$ on the line $l+1$.
Apply this $j\le 3$ times to the code of Theorem~\ref{fam} for the largest $l_0\equiv 1\pmod 4$ with $l_0\le l$; the resulting $k=(l_0^2-1)/2-j$ is in the range of the conjecture.
Lines $6,7,8$: Table~\ref{tab}.
\end{proof}
\begin{remark}[$l\equiv 3 \pmod 4$]
For these $l$ the type word of period $l$ is impossible at $2k+1=l^2$ (below), but inserting two extra $FF$ at antipodal places into
$2(l-1)$ periods gives a word of length $N=2l^2-2l+4=4k+2l+2$ with $k=(l-1)^2/2$. With the same triples $\{s,s+K,s+2K,s+3K+1\}$ at the
first $F$ of every period, a half-turn automorphism and antipodal dots, a short search (no SAT) completes it to a valid code for
$l=11,15,19,23,27$: $(n,k,N)=(81,50,224)$, $(155,98,424)$, $(253,162,688)$, $(375,242,1016)$, $(521,338,1408)$ (files \texttt{codes/dfam\_*}),
and the symmetric $(31,18)$ codes of length $88$ have this shape too. For $l\ge 15$ the completion is given by a closed form (\texttt{PROOF2.md}, Section 6b; \texttt{struct/l3/dfam2.py}):
the word is $(FF\,W_0^{\,l-1})^2$ with $W_0=FF.F.F\cdots F.$, the triples start at position $0$ and at the first letter of the periods $1,\dots,l-2$,
and the remaining $F$ positions are joined in pairs at distance $N/4+d$, $d\in\{1,2,3,4,-2\}$. It is verified from the definition for
$l=15,19,\dots,99$ and proved for $l\ge 43$ by the three steps of Theorem~\ref{fam} (with a cruder count in Step 1); a second AI instance
checked the proof. So $k_{\max}(l)\ge (l-1)^2/2$ for every odd $l\ge 7$ (with equality for $l=7$, and $k_{\max}(l)=(l^2-1)/2$ for $l\equiv 1\pmod 4$).
The proof is not yet written out in this note.
\end{remark}
For $l\equiv 3 \pmod 4$, $l\ge 7$, by contrast, no code of length $4k+2l+2$ exists at $2k+1=l^2$ (the argument for $(40,24)$ applies verbatim:
the period is forced, and the triple needs an $F$ at residue $s+\tfrac{l+1}{2}$, which is a dot residue because $\tfrac{l+1}{2}$ is even).

\section{The last spread on each line}\label{sec:last}
Let $k_{\max}(l)$ be the largest $k$ (of the parity opposite to $l$) for which a code of length $4k+2l+2$ exists on the line $l$.
Theorem~\ref{main} gives $k_{\max}(l)\le (l^2-1)/2$ for odd $l$ and $k_{\max}(l)\le (l^2-l-2)/2$ for even $l$ (for $l\equiv 2 \pmod 4$
this number is even, so $k_{\max}(l)\le (l^2-l-4)/2$ there). Theorem~\ref{fam} shows equality for $l\equiv 1\pmod 4$.
In this section $N=4k+2l+2$, so Proposition~\ref{shape} applies with $m=n-k$ (\eqref{H} holds in all cases below), and
$\lfloor (L+2m-k)/2\rfloor=\lfloor (L+l+1)/2\rfloor$.
\begin{theorem}\label{last}
{\rm(a)} If $l\equiv 0\pmod 4$, $l\ge 8$ and $k=(l^2-l-2)/2$, there is no circuit code of length $4k+2l+2$.\\
{\rm(b)} If $l\equiv 3\pmod 4$, $l\ge 11$ and $(l^2-l-2)/2\le k\le (l^2-1)/2$ ($k$ even), there is no circuit code of length $4k+2l+2$.\\
In both cases $K(n,k)=4k+2l$. Hence $k_{\max}(l)\le (l^2-l-6)/2$ for $l\equiv 0,3\pmod 4$.
\end{theorem}
Longer lengths are excluded by Lemmas~\ref{L1} and \ref{L2} as before. Codes of length $4k+2l+2$ at $k=(l^2-l-6)/2$ exist for
$l=7,11$ (found by the method of Section~\ref{sec:shape}) and, by explicit constructions (Theorem~\ref{fams} below),
for every $l\equiv 0\pmod 4$, $l\ge 8$, and every $l\equiv 3\pmod 4$, $l\ge 15$; for $l\equiv 2\pmod 4$
codes at $k=(l^2-l-4)/2$ exist for $l=6,10,\dots,30$ and $l=38$ (found by a satisfiability search inside a class of
reflection-symmetric codes; the largest is a $(1071,701)$ code of length $2882$). So
\[ k_{\max}(l)=\tfrac12(l^2-1),\quad \tfrac12(l^2-l-4),\quad \tfrac12(l^2-l-6),\quad \tfrac12(l^2-l-6)\qquad\text{for } l\equiv 1,2,3,0 \pmod 4 \]
holds for all $7\le l\le 30$; for $l\equiv 0,1,3\pmod 4$ it is a theorem for all $l$ (Theorems~\ref{fam}, \ref{last}, \ref{fams}), and for
$l\equiv 2\pmod 4$ the upper bound is proved while the lower bound is known for $l\le 30$, $l=38$ and conjectured in general.

Both parts use one more elementary fact. A window of length $K+2=k+3$ has $\odd\ge k+1$ and contains no coordinate three times, so it
contains at most one pair $(i,i^+)$:
\begin{equation}\label{Y2} \text{for every position } t:\qquad \text{not both } e_t\le 1 \text{ and } e_{t+1}=0. \end{equation}

\begin{proof}[Proof of (a)]
Here $K=l(l-1)/2$, $N=2(l^2-1)$, $L_0=2K$, there are $l(l+1)$ $F$ positions, and with $p=N-4K=2l-2$ we have $N=(l+1)p$ and
$K=(l/4)p\equiv 0\pmod p$. By \eqref{bal} every window of length $2K$ has $F(W)\le l^2/2$, which is also the average; so
$F(W)=l^2/2$ for all of them, the type word is invariant under the shift by $2K$, hence has period $\gcd(2K,N)=p$, with $l>p/2$
letters $F$ per period. So some two adjacent positions $s,s+1$ are both $F$ (an \emph{$FF$ pair at $s$}), and then there are $FF$ pairs
at $s+K$ and $s+2K$. For an $FF$ pair at $s$ the window $[s,s+2K+2)$ has $F(W)=l^2/2+2$, one more than the floor in \eqref{bal},
so it contains three consecutive occurrences $i,i^+,i^{++}$ of a coordinate, i.e.
\begin{equation}\label{C1} e_s+e_{s^+}\le 1\qquad\text{or}\qquad e_{s+1}=e_{(s+1)^+}=0 . \end{equation}
Suppose $(e_s,e_{s^+})=(0,0)$: the coordinate sits at $s,s+K,s+2K$. In \eqref{C1} for the pair at $s+K$ the second alternative
contradicts \eqref{Y2} ($e_{s+K}=0$), so $e_{s+2K}\le 1$; then the fourth gap of this coordinate has $e=p-e_{s+2K}$, so the first
alternative of \eqref{C1} fails for the pair at $s+2K$, and the second one gives $e_{s+2K+1}=0$, against \eqref{Y2} at $t=s+2K$.
The cases $(0,1)$, $(1,0)$ and $e_{s+1}=e_{(s+1)^+}=0$ are refuted in the same way by following the coordinate through the pairs at
$s+K$ and $s+2K$ (in the last case one ends with four gaps equal to $K$, whose sum is not $N$).
\end{proof}

\begin{proof}[Proof of (b)] Write $k=(l^2-1)/2-2j$, $0\le j\le (l+1)/4$; $j=0$ is the argument given for $(40,24)$. Let $j\ge1$.
Now $K=(l^2+1)/2-2j$ is odd, $N=4K+2l-2$, and there are $2K+2l$ $F$ positions.
Let $w$ be the walk on $\mathbb Z$ that steps $+1$ at an $F$ position and $-1$ at a dot (positions unrolled), and
$\Delta(a,b)=w(b)-w(a)$; one turn of the code raises $w$ by $2l+2$. By \eqref{bal}, $\Delta(t,t+L)\le l$ for odd $L$ and $\le l+1$ for
even $L$ in $[g_{\max},2K]$, and $\Delta(t,t+2K+1)\le l+2$, with equality only if $t,t+K,t+2K$ carry one coordinate. Taking
complements, $\Delta(t,t+B)\ge 1$ for all $B\ge 2l-1$; in particular each level $m$ has a last visit $u_m$, followed by two up-steps.
Put
\[ x_m=u_{m+l}-u_m-(2K-1). \]
Then $x_m\ge 0$ (as $\Delta(u_m,u_m+2K-1)\le l$), $x_m$ is even, $x_{m+2l+2}=x_m$ and $\sum_{m=1}^{2l+2}x_m=lN-(2l+2)(2K-1)=8j$.
If $x_m=0$ then $\Delta(u_m,u_m+2K+1)=l+2$, so $q=u_m+K$ is an $F$ position. Since $K$ is odd and $(l+1)/2$ is even,
$w(q)\ne m+(l+1)/2$: either $w(q)\le m+(l-1)/2$ (``low'') or $w(q)\ge m+(l+3)/2$ (``high''). Let $f_m=u_{m+1}-u_m-l$. Directly from
the definitions,
\[ f_m+\dots+f_{m+l-1}=x_m-4j,\qquad f_{m-2}+f_{m-1}=-(x_m+x_{m+l}) . \]
In the low case $u_{m+(l-1)/2}\ge q$, i.e.\ $f_m+\dots+f_{m+(l-3)/2}\ge (l+1)/2-2j$; pairing consecutive terms and using the two
identities this becomes: the sum of $x$ over a certain set $m+P_0$ of indices modulo $2l+2$ is at least $m_0=(l+1)/2+2j$. In the
high case one gets the same for a set $m+Q_0$ (here one uses that the level $m+(l+1)/2+2$ cannot be reached before
$u_{m-(l-1)/2}+2K+1$, which controls a possible later return of the walk to the level $m+(l+1)/2$). Since $l^2\equiv 1 \pmod{2l+2}$,
$\tau=lm$ is a coordinate on $\mathbb Z_{2l+2}$ in which $m\mapsto m+l$ is the shift by one, and in it the two sets are intervals:
with $a=(l+1)/2$ and $\xi(\tau)=x_m$,
\begin{equation}\label{zero} \xi(\tau)=0\ \Longrightarrow\ \xi[\tau+1,\tau+a]\ge m_0\quad\text{or}\quad \xi[\tau-a-1,\tau-1]\ge m_0 . \end{equation}
\emph{Case $1\le j\le (l-3)/4$.} The total mass $8j$ is less than $2m_0$, so the intervals chosen in \eqref{zero} for two zeros always
intersect; having at most $a+1$ points on a cycle of $4a$ points, they have a common point $\pi$, and every zero lies in
$[\pi-a,\pi-1]\cup[\pi+1,\pi+a+1]$. So $\xi$ has at most $2a+1$ zeros; but it has at most $4j$ nonzero values, and
$4a-4j>2a+1$. (This case uses no property of the code beyond its type word.)\\
\emph{Case $j=(l+1)/4$.} Now the mass is $4a=2m_0$, $K\equiv a+1 \pmod{4a}$, and the identities behind \eqref{zero} are exact:
$\xi[\tau+1,\tau+a]=2a+G$ with $G=u_{m+(l-1)/2}-q\ge 0$ even (low), $\xi[\tau-a,\tau-1]=2a-1-g$ with $g=u_{m+a}-q$ odd (high).
If $G=0$, $q$ is itself a last visit, of the level with coordinate $\tau+a+1$; if that level also has $x=0$, its triple
continues the triple of $m$, one coordinate sits at $u_m,u_m+K,u_m+2K,u_m+3K$, and \eqref{Y2} at $t=u_m+2K-1=u_{m+l}$ shows
$\xi(\tau+1)\ne 0$. If $g=-1$, then $q-1=u_{m+a}$ and \eqref{Y2} at $t=q-1$ shows $\xi(\tau-a)\ne0$. If $g\ge 1$ then $\xi(\tau-a-1)\ne 0$.
So every zero $\tau$ has an adjacent interval of at most $a+1$ points and mass $\ge 2a$ which either has mass $\ge 2a+2$ or comes with
a nonzero value of $\xi$ at a prescribed end point ($\tau+1$, $\tau+a+1$, $\tau-a$ or $\tau-a-1$). If these intervals pairwise intersect
one gets a contradiction as before. Otherwise two of them are disjoint, each has mass exactly $2a$, and $\xi$ is supported on two
clusters of spans $s_1,s_2\le a+1$ separated by two gaps of zeros of sizes $\gamma,\gamma'$. A zero in a gap can only use an interval that
contains a whole cluster and begins or ends at a nonzero point, or that reaches across a cluster and the other gap into the second
cluster; counting the possible positions gives
\[ \gamma\ \le\ \pos{a+1-s_1-\gamma'}+\pos{a+1-s_2-\gamma'}+3 \]
and the same with $\gamma,\gamma'$ exchanged, which is incompatible with $\gamma+\gamma'=4a-s_1-s_2$ for $a\ge 6$.
\end{proof}
\noindent The details of the index bookkeeping are in the accompanying file \texttt{PROOF2.md} (Sections 7--9); each of the three
arguments was re-derived by an independent AI instance, which also tested every intermediate statement on the $2{,}181$ type words
of $(87,54)$ and on $64$ real codes. For $l=11,15,19$ the statement (b) was first obtained by a SAT solver on the type word alone.

\begin{theorem}\label{fams}
For every $l\equiv 0\pmod 4$, $l\ge 8$, and every $l\equiv 3\pmod 4$, $l\ge 15$, there is an explicit circuit code of length $4k+2l+2$
at $k=(l^2-l-6)/2$; for every $l\equiv 2\pmod 4$, $l\ge 10$, there is one at $k=(l^2-l-8)/2$. Consequently
$k_{\max}(l)=(l^2-l-6)/2$ for all $l\equiv 0,3\pmod 4$, $l\ge 7$, and $(l^2-l-8)/2\le k_{\max}(l)\le (l^2-l-4)/2$ for $l\equiv 2\pmod 4$.
\end{theorem}
\noindent The codes have the automorphism $t\mapsto t+N/2$, their dots are antipodal, and the $F$ coordinates are given by an explicit
gap function $x\mapsto x+K+e(x)$ on a periodic type word with one shortened period per half (files \texttt{PROOF3.md}--\texttt{PROOF5.md}).
The proofs have three steps. (C) Following $x\mapsto x+K+e(x)$ row by row shows that every $F$ coordinate closes after four steps
(most rows lie in $2$-cycles $x\to x^+\to x+N/2$, a few in genuine $4$-cycles). (2--3) For windows longer than every gap,
$\odd(W)\ge k$ is the balance inequality \eqref{bal}; since the type word has no two adjacent dots, $2F(W)-L$ equals the number of
$FF$ pairs inside $W$ plus $0,\pm1$, the pairs are spaced by the period, and the only windows that exceed the floor are those
containing one pair too many, which contain one explicit triple at their first pair. (1) For windows of length $K+y$ shorter than
the largest gap, $\odd(W)=L-2\,ch(W)$ where $ch(W)$ counts the ``chords'' $[x,x+e(x)]$ inside the first $y$ positions $J$; as every
$F$ position of $J$ starts a chord, $ch(W)=F(J)-X(J)$ with $X(J)$ the number of chords leaving $J$ at its right end, and the claim
becomes $X(J)\ge\lfloor (2F(J)-y)/2\rfloor$, which is at most $2$ and is settled by looking at the few rows with small $e$ just before
the right end of $J$.


\section{Codes four longer than $4k+2l$}\label{sec:four}
Lemma~\ref{L1} with $L=2K$ also limits codes of length $4k+2l+e$ for larger even $e$: on an even line ($k$ odd) such a code needs
\[ k\ \le\ k_e(l)=\frac{2l^2+(3e-8)l+(e-4)^2-8}{2e}, \]
so $k_2=(l^2-l-2)/2$, $k_4=(l^2+2l-4)/4$, $k_6=(l^2+5l-2)/6$. For $e=4$ this is attained for every $l$ divisible by $4$.

\begin{theorem}\label{four}
Let $l\equiv 0\pmod 4$, $l\ge 8$, $k=(l^2+2l-4)/4$ and $n=(3k+l+1)/2$. Then $K(n,k)=4k+2l+4=l(l+4)$.
\end{theorem}
\noindent The first cases are $K(33,19)=96$, $K(68,41)=192$, $K(115,71)=320$.
\begin{proof}
\emph{Upper bound.} $k$ is odd and $k>3l-7$, so $N\le 5k+3$ by Lemma~\ref{L2}. For $N=4k+2l+e$, $4\le e\le k+3-2l$, Lemma~\ref{L1} with
$L=2K$ requires $f(e)\le 0$, where $2f(e)=(k+l-1+e)(4k+8-2l-e)-(4k+2l+e)(k+1)$. One finds $f(4)=4k+4-l^2-2l=0$ and
$2f(e)-2f(4)=(e-4)(2k+4-3l-e)$, which is positive for $6\le e\le k+3-2l$ because then $2k+4-3l-e\ge k+1-l>0$.

\emph{The code.} Put $K=k+1=l(l+2)/4$, $N=l(l+4)=4K+2l$; then $N/4=K+l/2$, $N/2=2K+l$, $K\equiv l/2\pmod l$ and $K$ is even.
Define $g$ on $\mathbb{Z}_N$ through $r=x \bmod l$:
$g=K$ for $r=1$; $g=K+2l-2$ for $r=2$; $g=K+1$ for $r\in\{l/2,l/2+1\}$; $g=N/4$ for $r\in E$, the even numbers of $[4,l/2-2]$ and the
odd numbers of $[l/2+3,l-1]$; and $g=N/2$ for the remaining $l/2-1$ residues. The map $x\mapsto x+g(x)$ permutes $\mathbb{Z}_N$; on residues it is
$1\to l/2+1\to 2\to l/2\to 1$ and the identity elsewhere. Its orbits (four elements with gaps $K,K+1,K+2l-2,K+1$; four equally
spaced elements; two antipodal elements) are the coordinates, $T_x$ is the orbit of $x$, and there are $(l+4)(3l-2)/8=n$ of them.

\emph{Validity.} All gaps are $\ge K$, and a window has the same odd coordinates as its complement, so we need $\odd(W)\ge k$ for
$W=[t,t+L)$, $L=K+y$, $0\le y\le K+l$. In such a window a twice-used coordinate occurs at most once, an equally spaced one at most
twice, and one of the first kind at most three times, three times only as $x,x^+,x^{++}$ with $x\equiv 1$ or $l/2$ (span $2K+1$).
Hence $\odd(W)=L-2(P-T)$, with $P$ the number of pairs $(x,x^+)$ in $W$ and $T$ the number of such triples, and the claim is
\begin{equation}\label{four*} P(t,y)-T(t,y)\ \le\ \lfloor (y+1)/2\rfloor . \end{equation}
For a set $S$ of residues write $|S\cap[a,b)|$ for the number of $x\in[a,b)$ with $x\bmod l\in S$. A pair starting at $x$ lies in $W$
iff $t\le x<t+L-g(x)$. Moving each pair to $x+g(x)-K$ gives, with $E'=E+l/2$ and
$M=\{0,1,l/2+1,l/2+2\}\cup E'$ (the odd numbers of $[1,l/2+1]$ and the even numbers of $[l/2+2,l]$),
\begin{align*} P(t,y)&=|M\cap[t,t+y)|-R(t,y),\\ R(t,y)&=[\,y\ge1,\ t\equiv l/2+1\text{ or }l/2+2\,]+|E'\cap[t,t+\min(y,l/2))|+|\{0\}\cap[t,t+\min(y,2l-2))|, \end{align*}
and $T(t,y)=|\{1,l/2\}\cap[t,t+y-K-1)|$.

(A1) $P(t,K)=K/2$ for every $t$: as $K=l\cdot l/4+l/2$, this reduces to $|\{0,1,l/2+1,l/2+2\}\cap[t,t+l/2)|=z(t)+d(t)$ with $z$ the
number of multiples of $l$ in $[t,t+2l-2)$ and $d$ the first term of $R$; both sides are $1,1,2,3,3,2$ for
$t\bmod l=1$; $2$; $3..l/2$; $l/2+1$; $l/2+2$; $l/2+3..l$.

(A2) $2l-2\le y\le K$: $R$ does not depend on $y$, and $M$ has no two adjacent non-members, so
$P(t,K)-P(t,y)=|M\cap[t+y,t+K)|\ge\lfloor (K-y)/2\rfloor$, i.e.\ $P(t,y)\le\lceil y/2\rceil$.

(A3) $y\le 2l-2$: now $P=|M'\cap[t,t+y)|-R'$ with $M'=M\setminus\{0\}$ and $R'$ the first two terms of $R$. Cut $[t,t+y)$ into
consecutive pairs. A pair holds two members of $M'$ only if it is $\{l/2+1,l/2+2\}$, none only if it is $\{l-1,l\}$ (mod $l$),
and these alternate; so $|M'\cap[t,t+y)|\le\lceil y/2\rceil+1$, with equality only if $t$ is odd, $1\le t\bmod l\le l/2+1$, and the
pair $\{l/2+1,l/2+2\}$ is inside. Then $R'\ge1$: either $t\equiv l/2+1$, or $t\in E'$, or $t\equiv 1$ and $t+2\in E'$.

(B) $y=K+d$, $1\le d\le l$: with $s=t+K$, $I=[s,s+d)$, $Z=\{0,l/2+1\}$ and $M''=M\setminus Z$, (A1) and $K\equiv l/2$ give
$P-T=K/2+|M''\cap I|+[\,s+d-1\in Z\,]$. $M''$ has no two adjacent members, so $|M''\cap I|\le\lceil d/2\rceil$; and if the last
position of $I$ is in $Z$, the last two positions of $I$ are non-members ($\{l-1,0\}$ or $\{l/2,l/2+1\}$), so
$|M''\cap I|\le\lceil d/2\rceil-1$. In both cases $P-T\le K/2+\lceil d/2\rceil=\lfloor (y+1)/2\rfloor$.
\end{proof}
\noindent Equality in Lemma~\ref{L1} at $L=2K$ means that every window of length $2K$ contains exactly $K/2$ pairs (this is (A1)) and
no twice-used coordinate twice; so the type word has period $\gcd(2K,N)$, which is $l$ here and $2l$ for $l\equiv 2\pmod 4$. The
family was read off from the solver code for $(33,19)$, which turned out to be invariant under $t\mapsto t+8$, and from an exhaustive
search of the codes invariant under $t\mapsto t+12$ for $(68,41)$. For $l=10$ the bound is attained as well, $K(49,29)=140$
(solver; type word of period $20$, no code invariant under $t\mapsto t+10$), but not for $l=6$ ($K(20,11)=58$).
The details are in \texttt{PROOF6.md}; an independent AI instance rebuilt the codes from the text for $l\le 64$ and checked every step.


\begin{theorem}\label{four2}
Let $l=8q+2$, $q\ge1$, $k=(l^2+2l-4)/4$ and $n=(3k+l+1)/2$. Then $K(n,k)=4k+2l+4=l(l+4)$.
\end{theorem}
\noindent The first cases are $K(49,29)=140$, $K(143,89)=396$, $K(285,181)=780$. The upper bound is as before. The code is invariant
under $t\mapsto t+2l$; one period of its type word is
\[ FF.\ \,FF.\ \,(F.)^{2q-2}\ \,FF.\ \,(F.)^{2q-1}\ \,FF.\ \,(F.)^{4q-1}, \]
and the excesses $e=g-K$ of the gaps of the $F$ positions are, block by block,
\[ (1,1)\ \ (1,\tfrac{3l}{2})\ \ (\tfrac l2-1)^{2q-2}\ \ (0,1)\ \ (\tfrac l2-1),(\tfrac{3l}{2})^{2q-2}\ \ (0,2l-2)\ \ 0^{2q-1},1,2l-2,1^{2q-2}. \]
The twice-used positions are paired at distance $N/2+1$ (residue $2$), $N/2$ (odd residues in $[5,4q+1]$) and $N/2-1$
(even residues in $[4q+4,l]$). The coordinates used four times have cyclic excess patterns $(0,1,1,2l-2)$, $(0,1,2l-2,1)$,
$(0,\frac{3l}2,\frac l2-1,1)$ and, $2q-2$ times per period, $(0,\frac l2-1,1,\frac{3l}2)$.
The proof (file \texttt{PROOF7.md}) follows the same four steps. Two of them are general: for any gap function invariant under a
shift $p$ with $2K\equiv 0\pmod p$, the number $X(t)$ of ``excess chords'' $[x,x+e(x)]$ with $x<t\le x+e(x)$ and the number of
members of $M=F-K$ in $[t,t+K)$ have the same increments in $t$, which gives (A1) by summing over a period; and (A2) only uses
that the type word has no two adjacent dots. (A3) and (B) become statements about intervals $J$ of residues: the number of
$F$ positions of $J$ whose preceding chord (resp.\ whose preceding triple does not) fit into $J$ is at most $\lceil |J|/2\rceil$,
and both follow by cutting $J$ into pairs of consecutive positions. An independent AI instance rebuilt the codes from the text
for $q\le 12$ and checked every step.

\begin{theorem}\label{none}
{\rm (a)} Let $l=8q+6$, $q\ge1$, $k=(l^2+2l-4)/4$. There is no spread-$k$ circuit code of length $4k+2l+4$ in dimension $(3k+l+1)/2$.
{\rm (b)} Let $l=8j+5$, $j\ge1$, $k=l(l+3)/4$ (the only case of equality in Lemma~\ref{L1} for $N=4k+2l+4$ on an odd line).
There is no spread-$k$ circuit code of length $4k+2l+4$ in dimension $(3k+l+1)/2$.
In both cases $K(n,k)\le 4k+2l+2$.
\end{theorem}
\noindent Together with Theorems~\ref{four} and \ref{four2}: for even $l\ge8$ and odd $l\ge13$, a code of length $4k+2l+4$ with
equality in Lemma~\ref{L1} exists if and only if $l\equiv 0,2,4\pmod 8$.
\begin{proof}[Sketch]
Let $L_0=2K$ ($k$ odd) or $2K-1$ ($k$ even). Equality in Lemma~\ref{L1} forces: every window of length $L_0$ contains the same
number $c_0$ of pairs $(x,x^+)$; twice-used coordinates have gaps $\ge L_0$; hence (shift the window by one) the type word has period
$L_0$, so period $p=\gcd(L_0,N)$. If $t,t+1$ are both four-times-used positions, the window $[t,t+L_0+2)$ contains $c_0+2$ pairs,
one more than $\odd\ge k$ allows, so it contains a triple $x,x^+,x^{++}$ with $x\ge t$, $x^{++}\le t+L_0+1$.
(a) Here $p=2l$ divides $K$, each period has $l+2>l$ such positions, so adjacent ones exist, say at residues $a,a+1$; there are
$N/2l=4q+5$ such pairs. The triple of a pair has gaps $(K,K)$, $(K,K+1)$ or $(K+1,K)$, so it stays in the residues $a,a+1$ and starts in that pair.
A coordinate has four positions with gap sum $N>4K+2$, so it contains at most two such triples, and two only if they are
consecutive and all four positions have residues $a,a+1$. Counting positions, $4q+5\le\sum_c s_c\le\frac12\sum_c b_c=4q+5$
($s_c$ triples, $b_c$ positions of $c$ at residues $a,a+1$), so every $s_c$ is $0$ or $2$: but $4q+5$ is odd.
(b) Here $p=(l+1)/2=4j+3$ with $2j+2$ four-times-used positions per period; comparing windows of lengths $L_0$ and $L_0-2$ shows that
no two twice-used positions are adjacent, so the type word is $FF.(F.)^{2j}$. The triple of the pair $FF$ at $t$ must have gaps $(K,K)$,
so $t+K$ is four-times-used; but $K\equiv 2j+2\pmod p$ is a twice-used position.
\end{proof}
\noindent (File \texttt{PROOF8.md}; checked by an independent AI instance, which also verified the forced properties on eight real
codes with equality.) The parity argument of (a) applies to any even excess $e$: it excludes every case of equality for $e=6$ and
$e=10$ --- in particular there is no $(30,17)$ code of length $90$, which a SAT solver confirms --- and one residue class of $l$ in four
for $e=8,12$. Most of the remaining cases are attained: see Section~\ref{sec:large}.

\section{Arbitrarily large excess}\label{sec:large}
The bound $k\le k_e(l)$ is attained for every excess $e$ divisible by $4$, on infinitely many lines each. Write $e=4j$.

\begin{theorem}\label{large}
Let $p\ge 10$ be even and $j\ge 1$, and put $k=jp(p+2)/4-1$, $n=j(3p-2)(p+4)/8$ (so $l=jp-2j+2$ and $k=k_{4j}(l)$).
If $p\equiv 0,2,4\pmod 8$ then
\[ K(n,k)\ =\ j\,p(p+4)\ =\ 4k+2l+4j ; \]
if $p\equiv 6\pmod 8$ there is no circuit code of that length, and $K(n,k)\le 4k+2l+4j-2$.
\end{theorem}
\noindent So $(n,k+1,K(n,k))=j\cdot\bigl((3p-2)(p+4)/8,\ p(p+2)/4,\ p(p+4)\bigr)$ for every multiplier $j$: every multiple of the
parameters of Theorems~\ref{four} and \ref{four2} is attained. Examples: $K(98,59)=280$ and $K(136,83)=384$ (excess $8$),
$K(272,167)=768$ (excess $16$), $K(544,335)=1536$ (excess $32$). The excess over Byrnes' value $4k+2l$ is therefore unbounded.
(For $p=8$ the statement holds for odd $j$ and is open for even $j$; the first case is $(66,39)$, $N=192$.)
For $e=4j$ with $j$ squarefree these are all the cases of equality in Lemma~\ref{L1} at $L=2K$.
In the proof below $p=4i$; the case $p\equiv 2\pmod 8$ is treated after it.
\begin{proof}[Sketch]
\emph{Upper bound} (all $p$). With $u=jp$, $K=k+1=jp(p+2)/4$, $e=4j$ one has $eK=u(u+e/2)$, which is $f(e)=0$ in the notation of the proof of
Theorem~\ref{four}; $k>3l-7$ and $2f(e')-2f(e)=(e'-e)(2k+8-3l-e-e')>0$ for $e<e'\le k+3-2l$ as there.

\emph{The code for $p=4i$, $i\ge 3$,} is invariant under $t\mapsto t+u$, $u=jp=4ji$; put $h=u/2$, $N=16ji(i+1)=4K+2u$, $X=2u-2$; then $K\equiv h\pmod u$
(residues are taken mod $u$).
A \emph{ladder} is a coordinate used four times whose gaps are $K+1,K+1,K,K+X$ in one of the three cyclic orders
$A=(1,1,0,X)$, $B=(1,1,X,0)$, $C=(1,0,1,X)$ of the excesses; started at residue $a$ it occupies the residues
$a,a+h+1,a+2,a+h+2$ (type $A$), $a,a+h+1,a+2,a+h$ (type $B$), $a,a+h+1,a+1,a+h+2$ (type $C$). For even $j$ take one ladder of
type $A$ at $a=0$, one of type $B$ at $a=4$ and $j-1$ of type $C$ at $a=8,13,18,\dots,5j-2$; for odd $j$ take the code of
Theorem~\ref{four} for $l=4i$ with $K$ and $N$ multiplied by $j$ and the long gap replaced by $N-3K-2$ (one ladder of type $C$,
period $4i$); $j$ ladders of type $C$ at $a=0,5,\dots,5j-5$ with period $4ji$ also work in every case computed.
Let $M_0$ be the set of the points $x+g(x)-K$ over the ladder residues $x$. In every maximal run of residues outside $M_0$ take
every second residue, shift this set by $h$, and give those residues the gap $N/4$; all remaining residues get the gap $N/2$.

\emph{Validity} follows the proof of Theorem~\ref{four}: $P(t,y)=|M\cap[t,t+y)|-R(t,y)$ where $M$ is the set of points
$x+g(x)-K$. (A1) $P(t,K)=K/2$ holds ladder by ladder, hence for any disjoint union of ladders of the three types. (A2) uses that $M$
has no two adjacent non-members. (B) uses that the adjacent pairs of members of $M$ are exactly the $2j$ pairs coming from the
ladders, each of which is the end of a triple of span $2K+1$. (A3), for $y\le X$: the points of $M$ whose pair can fit into a
window of length $K+y$, $y\le h$, form a sparse set $\Lambda$ ($3$ points per ladder) for which the pairing argument works
directly; for $h<y\le X$ one adds $|\Lambda\cap[t,t+h)|\le 2j+1$ and the bound $\lceil m/2\rceil+j/2$ for the number of members of
$M'$ in an interval of length $m$, and $2j+1+j/2\le h/2$ because $i\ge3$.
\end{proof}
\noindent\emph{The case $p=8q+2$.} Write the type word of Theorem~\ref{four2} (period $2p$) as $H_1H_2$ with
$H_1=(F.)^{4q-2}FF.FF.$ and $H_2=(F.)^{2q-2}FF.(F.)^{2q-1}FF.F.$ The code for the multiplier $j$ has type word $H_1^{\,j}H_2^{\,j}$,
period $2jp$: the four kinds of four-times-used coordinates of Theorem~\ref{four2} (excess patterns $(1,1,X,0)$, $(1,0,1,X)$,
$(0,b,a,1)$ and $(0,a,1,b)$ with $X=2jp-2$, $a=4q$, $b=2jp-1-a$) are repeated at the shifts $0,p,\dots,(j-1)p$, and the
twice-used positions are paired at distance $N/2$ or $N/2\pm1$ as before. The proof of Theorem~\ref{four2} goes through window by
window (\texttt{PROOF11.md}).

\noindent\emph{Where the bound is not attained.} The parity argument of Theorem~\ref{none}(a) needs only that the period
$\gcd(2K,N)$ of the type word divides $K$ and that the number of periods is odd. For $e\equiv 2\pmod 4$ this always holds: on
even lines the bound $k\le k_e(l)$ is \emph{never} attained when $e=6,10,14,\dots$ On odd lines ($L=2K-1$) the type word has an odd
period $p'$ with $2K\equiv 1$; reading it in the order $m\mapsto Km$, an adjacent pair of four-times-used positions forces a
triple with gaps $(K,K)$, which excludes the pattern $101$, and adjacent twice-used positions are excluded as before, which excludes
$0{*}0$; so the twice-used positions of a period come in pairs, while their number is odd whenever $e\equiv 0\pmod 4$. Hence on odd
lines the bound for $e=4,8,12,\dots$ is never attained either (\texttt{PROOF10.md}, sections 6 and 7).

\noindent\emph{A common generalisation.} The ladder designs do not need $u=4ji$: let $j\ge1$, $m$ odd and $p$ be such that
$4j\mid mp$ and $p/2\ge 5j+2$ ($j$ even) or $p/2\ge 5j+1$ ($j$ odd). Take the ladders above with period $p$ (for odd $j$: $j$ ladders of
type $C$ at $a=0,5,\dots,5j-5$), $u=mp$, $K=mp(p+2j)/(4j)$, $N=4K+2u$, $X=2u-2$, and the same forced filling. This is a spread-$(K-1)$
circuit code of length $N=4k+2l+4jm$, and $N=K(n,k)$ (\texttt{PROOF12.md}). For $m=1$ and even $j$ this is the code of the proof above; $j=1$
gives the codes of Theorem~\ref{four} with $K$ and $N$ multiplied by $m$; and it attains further cases of equality when $e/4$ is not squarefree,
e.g.\ $K(495,303)=1408$ ($e=36$, $u=96$).

\noindent The details are in \texttt{PROOF9.md}, \texttt{PROOF10.md}, \texttt{PROOF11.md} and \texttt{PROOF12.md}, each checked by an independent AI instance that rebuilt
the codes from the text. The proof template shows what a shift-invariant code at equality \emph{is}: a set of ladders whose
points do not touch, with a forced filling. Searching such designs directly takes milliseconds where the solver timed out; for
$e=16$ there are $46$ designs of this kind for $p=48$ and $1036$ for $p=64$, and none for $p=32$. Open: $p=8$ with $j$ even, the cases of equality for $e=4j$ with $j$ not squarefree and $j\nmid u$ (first $e=36$),
and, on odd lines, $e\equiv 2\pmod 4$, $e\ge 6$ (first $(70,42)$, $N=200$).

\section{Symmetric codes with a long bit run}\label{sec:sym}
Byrnes' conjecture compares $K(n,k)$ with $S(n,k,r)$, the maximum length of a spread-$k$ code that is \emph{symmetric}
($N=2H$ and $T_{i+H}=T_i$ for all $i$) and has a \emph{bit run} of length $r$ ($r$ cyclically consecutive distinct transitions).
Conjecture~4.2 of \cite{B2} is $S(n,k,k+l)=K(n,k)$; Byrnes proves $S(n,k,k+l)=4k+2l$ for $k\ge 2l+1$ ($k$ odd), $k\ge 2l-2$ ($k$ even).
The restriction on $k$ is not needed.

\begin{theorem}\label{symall}
Let $k\ge 1$, $a\ge 2$ and $2n\le a+2k+1$. A symmetric spread-$k$ circuit code in $Q_n$ with a bit run of length $a$ has length
at most $2a+2k$. In particular, for all $k,l\ge 1$ with $2n=3k+l+1$, $S(n,k,k+l)\le 4k+2l$; and $S(n,k,k+l)=4k+2l$ for all
$l\ge 2$, $k\ge l-1$. So Conjecture~4.2 holds at $(n,k)$ if and only if $K(n,k)=4k+2l$.
\end{theorem}
\begin{proof}
Shift the code so that the run is $A=\{0,\dots,a-1\}$, and let $f(p)=\odd([0,p))=d_H(v_0,v_p)$.
\emph{Two centres.} A coordinate of the run occurs once in $[0,a)$, so it is odd in exactly one of $[0,p)$ and $[a,p)$ ($p\ge a$);
any other coordinate is odd in at most both. Hence $\odd([0,p))+\odd([a,p))\le a+2(n-a)=2n-a\le 2k+1$.
\emph{Shell.} If $a+k\le p\le N-k$, both windows have length in $[k,N-k]$, so both counts are $\ge k$ by (F1), and therefore
$k\le f(p)\le k+1$; as $f(p)\equiv p\pmod 2$, $f$ alternates between $k$ and $k+1$ on this range.
\emph{End.} Suppose $H\ge a+k+1$. Then $p=H-1,H,H+1,H+2$ lie in the range. By symmetry $T_H=T_0$ and $T_{H+1}=T_1\ne T_0$,
and $T_{H-1}\ne T_0$ by (F2) (it is $K$-close to $T_H=T_0$; $N\ge 2k+2$).
If $f(H)=k+1$, then $f(H-1)=k$, and $T_0$ is odd in $[0,H)$ (else $f(H+1)=k+2$, as $T_H=T_0$); as $T_{H-1}\neq T_0$ it is odd in
$[0,H-1)$ and even in $[1,H-1)$, so $\odd([1,H-1))=f(H-1)-1=k-1$, contradicting (F1) for a window of length $H-2\in[k,N-k]$.
If $f(H)=k$, then $T_0$ is even in $[0,H)$ (else $f(H+1)=k-1$), $f(H+2)=k$, and $T_0$ is odd in $[0,H+2)$ (the transitions
$T_H=T_0$, $T_{H+1}\ne T_0$ add one occurrence); so $\odd([1,H+2))=k-1$ for a window of length $H+1\in[k,N-k]$, again a
contradiction. Hence $H\le a+k$. The lower bound $S\ge 4k+2l$ for $k\ge l-1$: Byrnes' sequence (eq.~(8) of \cite{B2}),
$(w_1w_{2A}w_{2B})^2$ with $w_1=1,\dots,k+l$, $w_{2A}=2,4,\dots,2(l-1)$, $w_{2B}=g_1,2l,g_2,2l+2,\dots,g_m,k+l-1$
($g_i=k+l+i$, $m=(k-l+1)/2$), is a spread-$k$ code of length $4k+2l$ for every $k\ge l-1$; we checked this by a counting
argument over windows (Theorem~37 in the companion file \texttt{PROOF17.md}; machine-checked, Section~\ref{sec:formal}).
\end{proof}

So everything in Conjecture~4.2 beyond Byrnes' own bound is the question whether $K(n,k)=4k+2l$: true in the range of
Theorem~\ref{main}, false at the points of Table~\ref{tab} and Sections~\ref{sec:last}--\ref{sec:large}, and on every line $l\ge 6$
at some $k\ge 2l+1$.

\paragraph{The block bound is Plotkin's.} Cutting a window of length $jK$ into $j$ blocks of length $K$ gives $j+1$ code vertices
pairwise at Hamming distance $\ge k$ (adjacent ones at distance $K$, the blocks being bit runs by (F2)); Plotkin's double count
(sum of the pairwise distances of $j+1$ binary words $\le n\lfloor (j+1)^2/4\rfloor$) gives, for $N\ge j(k+1)+k$,
\[ j(k+1)+\Bigl(\tfrac{j(j+1)}2-j\Bigr)k\ \le\ \Bigl\lfloor\tfrac{j+1}2\Bigr\rfloor\Bigl\lfloor\tfrac{j+2}2\Bigr\rfloor n , \]
with an extra $+1$ on each window of an odd number of blocks when $k$ is even (parity of $\odd$). For $j=4$ and $j=6$ this is
exactly Lemma~\ref{L2} and its six-block analogue ($N\le 7k+5$ for $k\ge 2l$, $k$ even, or $k\ge 2l-3$, $k$ odd), and for every
$n\le 2k-1$ it bounds $K(n,k)$ by $(2s+1)k+2s-1$ whenever $(s+1)(l+1)<(s-1)k+4$. This is weak in the tables
(e.g.\ $78$ against $K(15,8)=42$) and is very likely in \cite{Si} or \cite{Do}; we record it because Lemma~\ref{L2} is a case of it.

\section{Machine verification}\label{sec:formal}
The following are theorems of a Lean~4 + Mathlib development (about $11{,}300$ hand-written lines plus $9{,}500$ generated lines of value certificates; file \texttt{ccformal/FORMAL.md} lists each
statement), checked by the Lean kernel with no axioms beyond Lean's standard three (no \texttt{sorry}, no \texttt{native\_decide}).
Definitions: a code is a function $v:\mathbb N\to(\mathrm{Fin}\,n\to\mathrm{Bool})$ with Byrnes' definition; an independent audit
checked the definition against $57$ table entries and $194$ code files.
\begin{itemize}\setlength\itemsep{0pt}
\item Lemmas~\ref{L1}, \ref{L2}, Theorem~\ref{main}, the block bound above for every $j$, and the six-block lemma;
\item the obstructions of Sections~\ref{sec:shape}--\ref{sec:last} that settle the upper side of $k_{\max}(l)$ for every $l$
  (including the equality-point parity theorems), and the case $l=7$ of the last theorem of Section~\ref{sec:last};
\item the explicit families for $l\equiv1\pmod 4$ (all $l\ge 9$ and all larger excesses $4i+2$) and the excess-$4m$ family for
  $l\equiv 0\pmod 4$, $m$ odd, \emph{for all parameters} (via a general ``gap criterion'' for a transition word to be a code);
\item Byrnes' sequence for all $k\ge l-1$, Theorem~\ref{symall}, and hence: for every $l\ge 6$ there is $k\ge 2l+1$ with
  $S(n,k,k+l)=4k+2l<K(n,k)$, and $K-S$ is unbounded on the conjecture's range;
\item $k_{\max}(l)$ for every $l=6,\dots,31$ (a kernel-checked code at $k_{\max}(l)$ and $K=4k+2l$ for every larger $k$ on the line),
  and $94$ further exact values $K(n,k)$ via a verified checker.
\end{itemize}
Not formalized: the families for $l\equiv 0,2,3\pmod 4$ at general parameters (only instances), the SAT upper bounds of
Section~\ref{comp} (DRAT-checked instead), and one half of the equality-point theorem of Section~\ref{sec:four} ($u\ge 2e$).

\section{Counterexamples and computations}\label{comp}
\paragraph{Search.} The codes for $(22,12)$, $(23,13)$, $(25,14)$, $(28,16)$, $(30,17)$, $(31,18)$ were found by a SAT solver
(kissat) on a direct encoding of the definition for a fixed length and a prescribed longest run of distinct leading transitions
($k+4$). Deleting a coordinate that is used twice from a spread-$k$ code of length $N$ in $Q_n$ gives a spread-$(k-1)$ code of length
$N-2$ in $Q_{n-1}$ (Deimer \cite{De}; Theorem~4.2 of \cite{B1}); this maps $(n,k,l)\mapsto(n-1,k-1,l+1)$ and preserves the excess
over $4k+2l$, and gives $(29,16)$, $(36,20)$, $(37,21)$, $(39,22)$. The remaining codes were produced without SAT by a local ``pump''
$(n,k,N)\mapsto(n+3,k+2,N+8)$, which inserts one new coordinate at four places and two new twice-used coordinates into an existing
excess-$2$ code and keeps the result when it is valid. In every case the final code was checked from the definition.

\paragraph{Upper bounds by SAT.} For fixed $(n,k)$ split by the length $R$ of the longest run of distinct
transitions ($k+1\le R\le n$); a code of length $N$ whose longest run is $R$ contains a spread-$k$ open path
starting with that run, and unsatisfiability of ``path with $E$ edges starting with a maximal run $R$'' for all $R$ bounds $N$.
This proved $K(15,8)=42$, $K(18,10)=50$, $K(21,12)=58$, $K(24,14)=66$, $K(17,9)=50$, $K(20,11)=58$, $K(23,13)=66$, $K(26,15)=72$
(lower bounds: Byrnes' codes, our searches, and Table~\ref{tab}). For all eight, DRAT proofs of every unsatisfiable instance were produced with cadical and verified with drat-trim. The encoder, as opposed to the proofs, has only been validated empirically
(it reproduces all known optima and excludes none of several thousand known codes).

\paragraph{What is open.} Let $k_{\max}(l)$ be the largest $k$ with $K(n,k)>4k+2l$ on line $l$. Then $k_{\max}(5)=6$ and
$k_{\max}(l)=13,18,25,40,43,52,63,84,89,102$ for $l=6,\dots,15$; Section~\ref{sec:last} gives $k_{\max}(l)$ for all $l\not\equiv 2\pmod 4$,
and for $l\equiv 2\pmod 4$ the value $(l^2-l-4)/2$ is attained for every $l\le 38$ by solver codes (in a reflection-symmetric class),
with no formula and no proof for general $l$. Below $k_{\max}(l)$ the exact values are mostly open; for $k\le 3l-3$ (e.g.\
$K(31,18)\ge 88$, $K(22,12)\ge 64$) none of our upper-bound lemmas applies.

\paragraph{Caveats.} The proofs above were machine-tested (each inequality of Lemmas~\ref{L1} and \ref{L2} on about $19{,}000$ codes;
the arithmetic of Theorem~\ref{main} for $l\le 60$, $k\le 4000$) and read adversarially by a second AI instance, but by no human; later sections were added the same day; Section~\ref{sec:formal} lists what has since been checked by a proof assistant.
\paragraph{Prior work.} The prior-art search covered \cite{B1,B2} and both papers of Douglas \cite{Do} in full.
Douglas' first paper proves $K(n,k)=4k+2l$ on the lines $l=2$ (all odd $k$, with uniqueness of the code), $l=3$ (all even $k$)
and $l=4$ (odd $k\ge 9$), by a case analysis of transition sequences that contain a block of $k+2$, $k+3$ or $k+4$ distinct
transitions, and states without proof the bounds $4k+3l-1$ and $4k+3l+2$ for codes containing a block of $k+l$ distinct transitions;
it has no statement for general $l$. So on the lines $l\le 4$ Theorem~\ref{main} is Douglas' result (he covers more values of $k$
on the line $l=3$, we cover $k=7$ on the line $l=4$, where $K(13,7)=36$), and for $l\ge 5$ we found no earlier bound of this kind.
His second paper concerns $n$ large compared with $k$ (sphere-packing bounds) and does not touch the range $n<2k$.
The papers of Singleton and Deimer were not accessible to the author (Deimer's contraction is quoted from \cite{B2});
the block bound of Section~\ref{sec:sym} is Plotkin's bound applied to block boundaries.

\appendix
\section*{Appendix: a $(23,13)$ circuit code of length $66$}
Transition sequence (coordinates $0,\dots,22$; start at the zero vector):
\begin{quote}\small\raggedright
0 1 2 3 4 5 6 7 8 9 10 11 12 13 14 15 16 3 17 1 18 7 6 19 9 20 4 13 21 11 22 16 3
0 1 2 6 14 5 9 8 4 13 10 7 16 12 11 15 1 17 14 6 18 3 4 19 9 20 7 16 22 11 21 13 14
\end{quote}
Byrnes' Theorem~1.3 gives $64=4k+2l$ for $(n,k,l)=(23,13,6)$, the first case of his conjecture on the line $l=6$.

\begin{thebibliography}{9}
\bibitem{B1} K. M. Byrnes, \emph{The maximum length of circuit codes with long bit runs and a new characterization theorem},
Des. Codes Cryptogr. 87 (2019); arXiv:1809.08428.
\bibitem{B2} K. M. Byrnes, \emph{The maximum length and isomorphism of circuit codes with long bit runs}, arXiv:2008.04839.
\bibitem{De} K. Deimer, \emph{Some new bounds for the maximum length of circuit codes}, IEEE Trans. Inform. Theory 30 (1984).
\bibitem{Do} R. J. Douglas, \emph{Some results on the maximum length of circuits of spread $k$ in the $d$-cube}, J. Combin. Theory 6 (1969);
\emph{Upper bounds on the length of circuits of even spread in the $d$-cube}, J. Combin. Theory 7 (1969).
\bibitem{Si} R. C. Singleton, \emph{Generalized snake-in-the-box codes}, IEEE Trans. Electron. Comput. 15 (1966).
\end{thebibliography}
\end{document}
